
\subsection*{7.15. On fait bouger @@M0@@.}
Soit @@M1@@ un deuxième trait et
@@M2@@
un morphisme de traits, i.e. un morphisme local dominant (i.e. transformant point fermé en point fermé, point générique en point générique). On désigne par @@M3@@ (resp. @@M4@@) le point fermé (resp. le point générique) de @@M5@@, et on pose @@M6@@. Supposons @@M7@@ noethérien (p.ex. @@M8@@ de type fini sur @@M9@@). On a un diagramme commutatif
@@M10@@
à flèches obliques des morphismes plats. On peut donc prendre les homomorphismes images inverses pour @@M11@@ par toutes les flèches du diagramme (IV 2.12). Appliquant la définition 7.3, on trouve que le diagramme suivant
@@M12@@
est commutatif, où les flèches verticales sont les homomorphismes images inverses. On en déduit aussitôt un diagramme commutatif analogue, avec @@M13@@ au lieu de @@M14@@~; et lorsque @@M15@@ et @@M16@@ sont réguliers, on peut écrire encore @@M17@@ et @@M18@@ au lieu de @@M19@@ et @@M20@@~.

